\documentclass[12pt,reqno]{amsart}
\usepackage[T1]{fontenc}
\usepackage{lmodern}
\usepackage{amsmath,amssymb,microtype}
\usepackage[a4paper,left=32mm,right=32mm,top=27mm,bottom=27mm,headheight=14pt,headsep=15pt]{geometry}
\usepackage[hidelinks]{hyperref}
\hypersetup{pdftitle={Weak 3-randomness reflection does not imply Luzin's property (N)},
  pdfauthor={Zihao Guo},pdfsubject={Algorithmic randomness and computable real functions}}
\newtheorem{theorem}{Theorem}[section]
\newtheorem{lemma}[theorem]{Lemma}
\newcommand{\Cantor}{2^\omega}
\newcommand{\WThree}{\mathrm{W3R}}
\newcommand{\MLR}{\mathrm{MLR}}
\newcommand{\pre}{\mathbin{\upharpoonright}}
\newcommand{\Nprop}{\mathrm{(N)}}
\newcommand{\RR}{\mathbb R}
\newcommand{\mass}{\mu}
\newcommand{\lebesgue}{\lambda}
\newcommand{\eps}{\varepsilon}
\newcommand{\jump}{\emptyset'}
\newcommand{\djump}{\emptyset''}
\title[Weak 3-randomness reflection and Luzin's (N)]{Weak 3-randomness reflection does not imply\\ Luzin's property \texorpdfstring{$\Nprop$}{(N)}}
\author{Zihao Guo}
\date{November 21, 2026}
\subjclass[2020]{Primary 03D32. Secondary 03F60, 26A30.}
\keywords{Weak 3-randomness, randomness reflection, Luzin's property (N), computable real functions}
\begin{document}
\begin{abstract}
We construct a computable continuous real function which reflects both
weak 3-randomness and Martin--L\"of randomness relative to the second
Turing jump, but fails Luzin's property $\Nprop$. In particular, it maps
the set of reals failing weak 3-randomness into itself, answering a question
of Pauly, Westrick and Yu. The construction gives a localization theorem:
for every positive-measure $\Pi^0_1(\djump)$ class $A$ in Cantor space
there is a computable measure-preserving map that admits local effective
bounds on forward images off $A$, while sending a compact null subset
of $A$ to a set of measure at least $\mass(A)$.

\end{abstract}
\maketitle

\section{Introduction}

Which null sets suffice to detect that a computable real function fails
Luzin's property $\Nprop$? The property requires that every Lebesgue
null set have null image. A test of weak 3-randomness supplies a natural
candidate for detecting failure: its intersection is null and can be
specified by a sequence of open sets enumerable relative to the halting
problem. Pauly, Westrick and Yu ask whether a computable function can
take the union of all such test intersections to a null set and still
fail $\Nprop$ \cite[\S3.4, Open Question~29]{PWY}.

Write $\WThree$ for the weakly 3-random reals. A function $g$ reflects
a randomness class $\mathcal R$ if $g(x)\in\mathcal R$ implies
$x\in\mathcal R$ for every $x$. For $\mathcal R=\WThree$, reflection
says that $g[\RR\setminus\WThree]\subseteq\RR\setminus\WThree$.
The right-hand side is a countable union of null sets, so this inclusion
implies the null-image condition of Question~29. Westrick explicitly
asks whether such reflection implies $\Nprop$ \cite[slide~9]{W}.
The following example answers both questions.

\begin{theorem}\label{thm:main}
For every $0<\eps<1$ there are a computable continuous function
$g:\RR\to[0,1]$ and a compact set $B\subseteq[0,1]$ such that
\[
   \lebesgue(B)=0,\qquad \lebesgue(g[B])\ge 1-\eps .
\]
The same $g$ reflects both $\WThree$ and $\MLR^{\djump}$, where
$\MLR^{\djump}$ denotes Martin--L\"of randomness relative to $\djump$.
\end{theorem}

\noindent\textit{Formalization.} Theorems~1.1 and~2.1, Lemmas~2.2, 3.1, 4.1,
and~4.2, and the computable extension of Lemma~5.2 have been formalized in
Lean~4. The Lean proof of Theorem~1.1 establishes its full statement using
the test definitions in Section~2. For Lemma~5.1 the development proves the
homeomorphism and measure-scaling identity. It does not separately state
an algorithm computing finite prefixes of the inverse address from a fast
Cauchy name. The proof of Theorem~1.1 directly establishes the uniform
effective pullback of test layers along the embedding.
At \href{https://github.com/FAlepher/Weak_3_Randomness_and_Luzin_N/tree/4c52dff0973a7c2d9239bfe6c534c9066283b390}{commit \texttt{4c52dff} of the Lean repository},
\texttt{\#print axioms} reports only \texttt{propext},
\texttt{Classical.choice}, and \texttt{Quot.sound} for Theorem~1.1.

The $\MLR^{\djump}$ assertion also settles the case $\alpha=2$ of the
question about Martin--L\"of randomness relative to $\emptyset^{(\alpha)}$
in the paragraph following Open Question~28 of \cite{PWY}. The numbered
Question~28 itself concerns $\Delta^1_1$-randomness.

For comparison, $\Pi^1_1$-randomness reflection characterizes $\Nprop$
for computable functions \cite{PWY}. At lower levels the same authors
give a counterexample reflecting Martin--L\"of randomness, weak
2-randomness and Martin--L\"of randomness relative to $\jump$
\cite[Corollary~31]{PWY}. Their localization theorem makes the function
locally piecewise linear outside a $\Pi^0_1(\jump)$ class and places a
compact null witness of complexity $\Pi^0_1(\djump)$ inside it
\cite[Theorem~30]{PWY}. Such a witness is itself a null $\Pi^0_2(\jump)$
set. Its positive-measure image contains a weakly 3-random point, although
none of its inputs are weakly 3-random. This explains precisely why their
example fails to reflect $\WThree$. The descriptive-set-theoretic
construction of Holick\'y, Ponomarev, Zaj\'i\v cek and Zelen\'y
\cite{HPZZ} likewise gives an effective null witness in the instances
discussed immediately after Question~29 of \cite{PWY}.


For a computable real function the locus of local piecewise linearity
is $\Sigma^0_1(\jump)$ in the sense of \cite[Proposition~32]{PWY}.
Thus, if a computable function is locally piecewise linear at every
point that fails weak 3-randomness, it is locally piecewise linear
throughout $\RR$.\footnote{Any nonempty complementary
$\Pi^0_1(\jump)$ class meets a compact
interval with rational endpoints. A path search relative to $\djump$
in that interval yields a $\djump$-computable point.
Its singleton is a null $\Pi^0_2(\jump)$ set, so the point fails
weak 3-randomness.}

Our construction moves the prescribed closed class up one jump. Its
complement need only admit local bounds for the cost of covering forward
images. This change allows a computable map to act on a cylinder whenever
finite evidence supports it, even if that evidence eventually stops.
Inside the closed class the map retains enough of a second coordinate
to produce a compact null witness by taking a section. The closed class
can be chosen to consist entirely of $\MLR^{\djump}$ inputs.

\section{Localization and the choice of a random class}\label{sec:application}

Let $X=\Cantor$ with fair-coin measure $\mass$. Write $[\sigma]$ for the
cylinder with prefix $\sigma$, $x\pre k$ for the first $k$ bits of $x$, and
$p\preceq v$ for the prefix relation. Juxtaposition denotes concatenation.
For an oracle $Z$, a $Z$-c.e.\ open set is the union of a $Z$-enumerable
list of cylinders. A $\Pi^0_1(Z)$ class is the complement of such an open
set. All computability assertions about functions are in the usual Type-2
sense.

On $X$ we use the test characterization of $\WThree$ appearing in
\cite[Open Question~29]{PWY}. A point lies outside $\WThree$ if it belongs
to a null intersection $\bigcap_n U_n$ of uniformly $\jump$-c.e.\ open
sets. The $U_n$ can be taken decreasing by forming finite intersections.
A Martin--L\"of test relative to $Z$ consists of uniformly $Z$-c.e.\ open
sets $U_n$ with $\mass(U_n)\le 2^{-n}$. A point is in $\MLR^Z$ if it
avoids the intersection of each such test. On $\RR$ we use the corresponding
rational open intervals and Lebesgue measure. See also \cite{DH}.

To relate the weak test formulation to the usual effective
$\Pi^0_3$ null sets on $X$, observe that a $\jump$-c.e.\ open is an
effective $\Sigma^0_2$ set. Conversely, for an effective $\Sigma^0_2$
set $E=\bigcup_j P_j$ given as a union of effectively closed sets,
$\jump$ can cover each $P_j$ by a clopen set of measure less than
$\mass(P_j)+2^{-n-j-1}$. Indeed, $\jump$ computes $\mass(P_j)$ from
its decreasing clopen approximations. Taking the union gives a
$\jump$-c.e.\ open superset of $E$ of excess measure at most $2^{-n}$.
Applying this uniformly to the $n$th level of a decreasing effective
$\Sigma^0_2$ presentation of a null $\Pi^0_3$ set covers it by a null
$\Pi^0_2(\jump)$ set.

A useful model for failure of $\Nprop$ is the first-coordinate projection
$\pi:X^2\to X$. Each slice $X\times\{t\}$ has product measure zero and
projects onto $X$. Interleaving bits identifies $X^2$ computably and
measure-preservingly with $X$. The unrestricted projection, however, can
send a nonrandom pair $(y,0^\omega)$ to a random $y$. We shall realize the
slice argument on a closed class $A$ of random inputs. Off $A$ we need a
different estimate: a small neighborhood of a nonrandom input must have
an image that can be covered effectively and at bounded cost.

\begin{theorem}[Localization]\label{thm:local}
Suppose $A\subseteq X$ is a positive-measure $\Pi^0_1(\djump)$ class. There
is a computable $F:X\to X$ with $F_*\mass=\mass$ and the following properties.
\begin{enumerate}
\item\label{it:outside} Each $x\in X\setminus A$ belongs to a cylinder
  $I\subseteq X\setminus A$ with a constant $c<\infty$ such that, for
  every oracle $Z$, from an index for a $Z$-c.e.\ open $U$ one can uniformly
  obtain an index for a $Z$-c.e.\ open $V$ with
  \[
       F[U\cap I]\subseteq V,\qquad \mass(V)\le c\,\mass(U\cap I).
  \]
\item\label{it:inside} There is a compact $D\subseteq A$ such that
  $\mass(D)=0$ and $\mass(F[D])\ge\mass(A)$.
\end{enumerate}
\end{theorem}

The effective content of part~\eqref{it:outside} is local: for the
chosen $I$ there is one procedure on open-set indices, valid relative
to every oracle. Selecting $I$ from $x$ is an existence assertion.
Part~\eqref{it:inside} asserts compactness in the ordinary topology.

We first deduce the randomness assertion on Cantor space from this theorem.
The constant in part~\eqref{it:outside} may depend on the cylinder. Once a
cylinder containing a particular input has been fixed, the same constant
applies to every layer of a test for that input.

\begin{lemma}\label{lem:reflection}
For $0<\eps<1$ there is a class $A$ as in Theorem~\ref{thm:local} with
$\mass(A)\ge1-\eps$ for which $F$ reflects both $\WThree$ and $\MLR^{\djump}$.
\end{lemma}
\begin{proof}
Let $G$ be a layer of a universal $\djump$-Martin--L\"of test of measure
at most $\eps$ and put $A=X\setminus G$. Then $A\subseteq\MLR^{\djump}
\subseteq\WThree$. For the second inclusion, the measures of uniformly
$\jump$-c.e.\ open sets in Cantor space are uniformly left-c.e.\ in
$\jump$, hence computable in $\djump$. From any decreasing sequence of
such sets with null intersection, $\djump$ can select a subsequence with
measures less than $2^{-n}$. It is a test covering that intersection.

If $x\notin\WThree$, take a decreasing test $(U_n)$ with $x\in\bigcap_n U_n$.
As $x\in G$, part~\eqref{it:outside} of Theorem~\ref{thm:local} gives a cylinder $I$ containing $x$, a constant $c$, and
uniformly $\jump$-c.e.\ open $V_n$ containing $F[U_n\cap I]$, with
$\mass(V_n)\le c\mass(U_n)\to0$. Thus $F(x)$ belongs to the null
intersection $\bigcap_n V_n$ and is outside $\WThree$.

If $x\notin\MLR^{\djump}$, fix $I,c$ as above (since $x\in G$) and an integer $d\ge0$ with $2^d\ge c$. Take a
$\djump$-Martin--L\"of test $(U_n)$ covering $x$. Apply the same cover procedure
to $U_{n+d}$. The resulting $V_n$ contain $F(x)$ and have measure at
most $c\,2^{-(n+d)}\le2^{-n}$. This gives a $\djump$-Martin--L\"of test
covering $F(x)$.
\end{proof}

The localization theorem is the remaining Cantor-space assertion. Its proof occupies
Sections~\ref{sec:code} and~\ref{sec:image}. Section~\ref{sec:real}
then transports the resulting Cantor-space example to the real line.

\section{Coding with two finite strings}\label{sec:code}
The construction reads input bits in blocks and distributes them between
an output string and a temporary string. Bits already output are permanent.
When an input falls in a selected cylinder, the procedure moves most of
its temporary bits to the output. Repeated selections will give the local
cover estimate. On $A$, each fixed initial segment of the temporary string
eventually survives every selection, providing the second coordinate for
the slice argument.

\subsection{A computable schedule}
The program computing $F$ must decide what to do with the next input
bits using only finite evidence. An enumeration relative to $\djump$
cannot in general be run by this program. We therefore give each
candidate cylinder repeated opportunities to act. An input outside $A$
needs one cylinder that acts at arbitrarily large depths: this will let
us cover images of arbitrarily small neighborhoods. For an input in $A$,
we need only ensure that each fixed initial segment of its stored string
is disturbed finitely often. These two requirements suggest testing
candidate witnesses repeatedly and allowing a finite number of mistaken
selections for each unsuccessful witness.

Put $G=X\setminus A$. A $\djump$-c.e.\ list of cylinders for $G$ has a
$\Sigma^0_3$ index predicate on finite strings. By the arithmetical normal
form, for some computable $R$ a string $\sigma$ occurs on the list exactly
when
\[
        \exists e\ \forall n\ \exists t\ R(\sigma,e,n,t).
\]
Give each pair $(\sigma,e)$ a distinct computably decodable positive
index $i\ge|\sigma|$, using a computable coding with decidable range. Put
$I_i=[\sigma]$. Unused indices may be assigned an arbitrary cylinder and
are never selected. At computably reserved stages visit each used $i$
infinitely often, with at most one visit per stage $s$ and $i\le s$.
For example, reserve stages $2^{i-1}(2m+1)$ for $i$, $m\ge0$. At a visit,
compute the largest $\ell\le s$ such that for all $n<\ell$ some $t\le s$
satisfies $R(\sigma,e,n,t)$. Select $i$ whenever this length exceeds its
previous record. At other stages the schedule is idle. If $S$ is the set
of indices selected infinitely often, this gives
\begin{equation}\label{eq:schedule}
             G=\bigcup_{i\in S} I_i .
\end{equation}
Indeed, the recorded lengths are unbounded exactly for the pairs satisfying
the displayed $\forall n\,\exists t$ condition. All stages of the schedule
are computable. Only the description of their limiting behavior uses $S$.

\subsection{A finite bijection}
At the beginning of stage $s\ge1$ we have read $N_s=h_s+s$ bits and encoded
them as a pair $(u_s,v_s)$ of lengths $(h_s,s)$. The string $u_s$ is an
output prefix. The string $v_s$ stores the other bits. Initially $h_1=1$
and $(u_1,v_1)$ consists of the first and second input bits, respectively.
We maintain a bijection
\begin{equation}\label{eq:bijection}
   2^{N_s}\longrightarrow 2^{h_s}\times 2^s
\end{equation}
with the additional property that for each $j\le s$ we can recover
$x\pre j$ from $(u_s,v_s\pre j)$.

To derive the update, consider the cost of covering the image of a
single cylinder of depth $N_s$. A cover by one output cylinder of depth
$N_s-i$ costs $2^i$ times its input measure, independently of $s$.
To obtain this bound with a single output cylinder, it suffices to fix
a further $s-i$ output bits, since $(N_s-i)-h_s=s-i$. The stored
string supplies exactly that many bits once we reserve the first $i$.
Those reserved bits, together with $u_s$, recover the input prefix on
which membership in $I_i$ depends. On a hit we can thus release the
stored tail to the output and replace it with fresh bits. On a miss
we can send the fresh bits directly to the output, leaving the entire
stored string available at the next stage. Both choices have the same
two string lengths, so the branch need not be signalled by an extra bit.

If the schedule selects $i$ at stage $s$, write $u=u_s$ and $v_s=ab$
with $|a|=i$. Read a fresh block $c$ of length $s-i$. The intermediate
state is given by
\begin{equation}\label{eq:update}
(u,ab,c)\longmapsto
\begin{cases}
 (ub,ac),& x\in I_i,\\
 (uc,ab),& x\notin I_i .
\end{cases}
\end{equation}
At an idle stage keep the pair unchanged and read no block $c$. In either
case finish the stage by reading two fresh bits, appending the first to the
output and the second to the stored string. Thus $h_{s+1}=h_s+s-i+1$ at
a selection and $h_{s+1}=h_s+1$ at an idle stage.

\smallskip
\noindent\emph{A finite example.}
Suppose the first two stages are idle and stage $3$ selects $i=1$
with $I_i=[1]$. At this point $u$ holds the odd positions of the
six input bits and $v$ the even positions. Take the fresh block $c=11$.
The table ends just before the two final bits are appended.
\[
\begin{array}{c|ccc|cc}
 \text{old input} & u & a & b & \text{new output} & \text{new store}\\ \hline
 100110 &101&0&10&101\,10&0\,11\\
 000110 &001&0&10&001\,11&0\,10
\end{array}
\]
In the first row all four choices of $c$ give the same five output
bits. This is the promised cover of a six-bit input cylinder at cost
two. In the second row all four choices preserve the stored string
$010$. This is the feature used later to define a second coordinate on
$A$. Both rows represent selections of the same index.

The prefix $a$ is retained in both branches for a reason. Given $u$ and $a$,
the recovery property decides whether the input belongs to $I_i$. Thus the
new state records enough information to undo the exchange, even though it
contains no separate bit recording which branch was taken.

\begin{lemma}\label{lem:finite}
At each stage the encoding is a computable bijection with the stated
prefix recovery property. Consequently, for $|p|\le h_s$ and $|q|\le s$,
\begin{equation}\label{eq:product}
   \mass\{x:p\preceq u_s(x),\ q\preceq v_s(x)\}
       =2^{-|p|-|q|}.
\end{equation}
\end{lemma}
\begin{proof}
At stage~1 the two-bit code is the identity and its output bit
already recovers the first input bit. Suppose the claim holds at stage $s$.
In either branch of \eqref{eq:update}, the first $h_s$
output bits are $u$ and the first $i$ stored bits are $a$. The induction
hypothesis recovers $x\pre i$ from $(u,a)$. Comparing it with $\sigma_i$,
where $I_i=[\sigma_i]$ and $|\sigma_i|\le i$, determines the branch.
To see surjectivity as well as injectivity, take an arbitrary candidate
intermediate state $(uw,az)$, with $|u|=h_s$, $|a|=i$ and
$|w|=|z|=s-i$. In the first branch set $(b,c)=(w,z)$, and in the
second set $(b,c)=(z,w)$. Invert the old code on $(u,ab)$. The recovered
input has the prefix decoded from $(u,a)$ in either case, so it takes
exactly the branch used to reconstruct it. The two final bits can be
stripped off and restored independently.

For $j\le i$, the first $j$ stored bits survive the transition. For
$i<j\le s$, knowledge of the new output and the first $j$ stored bits
again identifies the branch. The first $j-i$ bits of $b$ are read from
the output in the first branch and from the stored string in the second.
This recovers $(u,v_s\pre j)$, hence $x\pre j$. At an idle stage these
prefixes are unchanged. For $j=s+1$ the full new state is available, so
bijectivity recovers the entire input read so far. This completes the
induction. A uniform input of length $N_s$ gives a uniform pair in
\eqref{eq:bijection}, which proves \eqref{eq:product}.
\end{proof}

Set $F(x)=\bigcup_s u_s(x)$. The output extends at every stage, and
$h_s\ge s$. For a requested output length $n$, a computable search for
$h_s\ge n$ gives an input bound $N_s$. Thus $F$ is computable and
continuous. Taking $q$ empty in \eqref{eq:product} shows $F_*\mass=\mass$.

\section{Forward images and the compact witness}\label{sec:image}

\begin{lemma}\label{lem:covers}
For each $i\in S$, the cylinder $I_i$ satisfies
Theorem~\ref{thm:local}\,\eqref{it:outside} with $c=2^i$.
\end{lemma}
\begin{proof}
Let $D_0\subseteq I_i$ be a cylinder. Wait for a stage $s$ at which $i$
is selected and $N_s$ is at least the defining length of $D_0$. For each
$\tau$ of length $N_s$ with $[\tau]\subseteq D_0$, the old pair $(u,ab)$
is fixed on $[\tau]$. Its points take the first branch of \eqref{eq:update},
so $F([\tau])\subseteq[ub]$. The bits of $b$ were read before stage $s$,
and
\[
       |ub|=h_s+s-i=N_s-i,\qquad \mass([ub])=2^i\mass([\tau]).
\]
These output cylinders cover $F[D_0]$ with total cost at most
$2^i\mass(D_0)$. The algorithm waits for a sufficiently late
selection of $i$. Its termination uses $i\in S$, while the instructions
for the search are computable from $i$ alone.

From an enumeration of $U\cap I_i$ enumerate disjoint cylinders for its
union, by taking the difference of each new cylinder with the previous
finite union. Run the finite-cover procedure on each resulting cylinder
in parallel and enumerate all outputs into $V$. The total cost bound
sums to $2^i\mass(U\cap I_i)$. The procedure is uniform in the enumeration
and the oracle $Z$. Equation~\eqref{eq:schedule} supplies a suitable $i$
for each $x\notin A$.
\end{proof}

The same cylinder estimate gives an outer-measure bound for every
$E\subseteq I_i$, $i\in S$. Cover $E$ by countably many cylinders in
$I_i$ and take the infimum of the resulting costs to obtain
\[
                \mass^*(F[E])\le 2^i\mass^*(E).
\]
Equation~\eqref{eq:schedule} now shows that $F$ sends every null subset
of $G$ to a null set.

\begin{lemma}\label{lem:witness}
There is a compact null $D\subseteq A$ with $\mass(F[D])\ge\mass(A)$.
\end{lemma}
\begin{proof}
Fix $k\ge1$. From stage $s\ge k$ onward, the first $k$ stored bits can
change at $x$ only when an $i<k$ is selected and $x\in I_i$. Every
$i\in S$ has $I_i\cap A=\varnothing$ by \eqref{eq:schedule}. Each of
the finitely many other indices $i<k$ acts only finitely often. Hence
there is a single $s_k\ge k$ after which these $k$ bits are constant at
every point of $A$. They define a map $T:A\to X$ with
\[
       T(x)\pre k = v_s(x)\pre k \quad(x\in A,\ s\ge s_k).
\]
This is a continuous map since each finite-stage reading is continuous.
The computation of $F$ only uses the schedule and requires no $s_k$.

Consider the compact set $H=(F,T)[A]$. A fixed second coordinate
will give a null input fiber. To show that some such fiber has a large
image we first estimate the product measure of $H$. For $n\ge1$ let
$\mathcal R_n$ consist of the rectangles $[p]\times[q]$ meeting $H$
with $|p|=|q|=n$, and put $H_n=\bigcup\mathcal R_n$.

\smallskip
\noindent\textit{Claim.} For $s\ge s_n$ let $E_{n,s}$ consist of all
$x\in X$ whose two stage-$s$ readings begin with a pair $(p,q)$
indexing a rectangle in $\mathcal R_n$. Then
\[
 A\subseteq E_{n,s},\qquad
 \mass(E_{n,s})=(\mass\times\mass)(H_n).
\]
Indeed, at each $x\in A$ the two prefixes read at stage $s$ coincide
with the prefixes of $F(x)$ and $T(x)$, so their rectangle meets $H$.
For the equality fix one rectangle. Its prefixes specify $n$ of the
$h_s$ output bits and $n$ of the $s$ stored bits. They leave exactly
$2^{h_s-n}2^{s-n}=2^{N_s-2n}$ complete states available. Each has
exactly one input word of length $N_s$ by Lemma~\ref{lem:finite}.
Consequently the inputs for this rectangle occupy measure
$2^{N_s-2n}2^{-N_s}=2^{-2n}$, its product measure. Distinct rectangles
give disjoint input sets. Summing over $\mathcal R_n$ proves the claim.

This count uses the bijection on the whole input space. Restriction to
$A$ enters only through the inclusion in the claim. Since $H$ is closed,
every point outside it has a prefix rectangle disjoint from it. Hence
$H_n\downarrow H$, and continuity of the finite product measure gives
\begin{equation}\label{eq:mass}
 (\mass\times\mass)(H)=\lim_n(\mass\times\mass)(H_n)\ge\mass(A).
\end{equation}

Fix $t\in X$. In the input space the fiber $D_t=T^{-1}(\{t\})$ is
compact, since $T$ is continuous on the compact set $A$. Each $x\in
D_t$ must already display $t\pre n$ in the first $n$ stored positions
at stage $s_n$. By \eqref{eq:product} this condition on the whole input
space has measure $2^{-n}$. Letting $n$ increase gives $\mass(D_t)=0$
for every $t$.

In the product space the same choice of $t$ gives the section
$H^t=\{y:(y,t)\in H\}$. By definition $(y,t)\in H$ exactly when there
is an $x\in A$ with $T(x)=t$ and $F(x)=y$. Thus $H^t=F[D_t]$.
Applying Fubini's theorem to the indicator of the compact set $H$ yields
\[
  \int_X \mass(F[D_t])\,d\mass(t)=(\mass\times\mass)(H)\ge\mass(A).
\]
Some $t$ has $\mass(F[D_t])\ge\mass(A)$: otherwise the positive
difference $\mass(A)-\mass(F[D_t])$ would have positive integral.
This $D_t$ is the required witness.
\end{proof}
Lemmas~\ref{lem:covers} and \ref{lem:witness} complete the proof of
Theorem~\ref{thm:local}.

\section{A computable function on the real line}\label{sec:real}

The two changes of space serve different purposes. An embedding into a
positive-measure Cantor set on the input side keeps real null tests null
when pulled back to $X$. Binary evaluation on the output side preserves
the lower bound for the image measure. Affine interpolation extends the
map across the gaps.

\begin{lemma}[Effective embedding and measure scaling]\label{lem:embed}
There is a positive-measure Cantor set $C\subseteq[0,1]$ and a
computable homeomorphism $h:X\to C$ with $h_*\mass=2\lebesgue\pre C$.
\end{lemma}
\begin{proof}
Start with $[0,1]$ and
use two closed subintervals abutting opposite ends of each parent.
Choose the common length at depth $n$ to be
\[
            L_n=2^{-n-1}(1+2^{-n}).
\]
Since $L_0=1$ and $2L_{n+1}<L_n$, the two children are disjoint and their
endpoints are rational. Let $C$ be the intersection of all levels.
Following the left and right children defines a computable homeomorphism
$h:X\to C$. If $|\sigma|=n$, then
$\lebesgue(h([\sigma]))=\lim_{m\to\infty}2^{m-n}L_m=2^{-n-1}$. Thus,
with $b:X\to[0,1]$ the usual binary expansion map,
\begin{equation}\label{eq:push}
       h_*\mass=2\,\lebesgue\pre C,\qquad b_*\mass=\lebesgue\pre[0,1].
\end{equation}

Indeed, the branch measures determine the restricted measure on $C$,
and each dyadic interval has the corresponding binary cylinder as its
preimage up to endpoints of measure zero. For the inverse of $h$ on
$C$, at each depth the finitely many rational closed intervals are
separated by positive gaps, so the unique interval containing the
given point can be identified from a sufficiently precise approximation.
\end{proof}

\begin{lemma}[Computable extension]\label{lem:extend}
For the map $F$ constructed in Section~\ref{sec:code}, the prescribed
values $bFh^{-1}$ on $C$ extend to a computable continuous
$g:\RR\to[0,1]$ by affine interpolation on the gaps and constant
extension on the exterior rays.
\end{lemma}
\begin{proof}

Define $g(h(x))=b(F(x))$ on $C$. Interpolate affinely on each bounded
complementary interval in $[0,1]$ and extend constantly on each outside
ray. To see that $g$ is computable, for any $k$ compute a level $m$ so
that $F$ has the same first $k$ output bits throughout each input
cylinder of length $m$. The oscillation of $bF$ on such a cylinder is
at most $2^{-k}$. At each endpoint of the level-$m$ intervals the
prescribed value is computable, since the corresponding path has a
constant tail. Linear interpolation between consecutive endpoints agrees
with $g$ on the earlier gaps. Within each retained interval both
interpolations take values in the closed convex hull of its values on
$C$, an interval of length at most $2^{-k}$. Approximating the endpoint
values by rationals to within $2^{-k}$ gives a rational polygonal
approximation with uniform error at most $2^{1-k}$. This also controls
the extensions on the rays.

\end{proof}

\begin{proof}[Proof of Theorem~\ref{thm:main}]
Choose $A,F$ by Lemma~\ref{lem:reflection}, take $D$ from
Lemma~\ref{lem:witness}, and let $g$ be the extension of
Lemma~\ref{lem:extend}. For a point in $C$ we follow a failed test
through three spaces. We pull a real test back along $h$, apply
reflection for $F$, and then cover the binary image of a failed test
by real intervals.

We verify reflection first at $h(x)\in C$. The inverse images under $h$
of uniformly $Z$-c.e.\ open sets are uniformly $Z$-c.e.\ open, and
\eqref{eq:push} gives $\mass(h^{-1}(U))\le2\lebesgue(U)$. If $h(x)$
fails weak 3-randomness, pulling back its null test shows that $x$ fails
as well. For a Martin--L\"of test relative to $\djump$, pull back the
$(n+1)$st layer to obtain the $n$th layer on $X$. In either case
Lemma~\ref{lem:reflection} transfers failure of the test notion to $F(x)$.

For a $Z$-c.e.\ open $U\subseteq X$, enumerate the interiors of the
dyadic intervals $b([\sigma])$ for a disjoint cylinder enumeration of
$U$. Their union $W$ has $\lebesgue(W)\le\mass(U)$ and contains every
nondyadic point of $b[U]$. Applying this operation to test layers
transfers nonrandomness of $z$ to $b(z)$ when $b(z)$ is nondyadic.
For a weak test on $X$ the layers can be chosen decreasing with measures
tending to zero, so the resulting intersection is again null. Dyadic
outputs are computable and fail both randomness notions. With $z=F(x)$
this proves the asserted reflection on $C$.

On a complementary interval, $g$ agrees with a computable affine map
$r\mapsto ar+d$. For $a\ne0$ its action sends effective opens to effective
opens, commutes with intersections by injectivity, and multiplies Lebesgue
measure by $|a|$. A fixed shift of the layer index absorbs this factor in
a Martin--L\"of test. If $a=0$ the value is computable, as are the
constants on the outside rays. This proves reflection on all of $\RR$.

The case split on the slope is used only to prove reflection at a
given point. The polygonal approximation computes $g$ uniformly across
all the gaps without deciding whether any slope is zero.

Finally put $B=h[D]$. It is compact and \eqref{eq:push} gives
$\lebesgue(B)=\mass(D)/2=0$. Since $F[D]$ is compact, the same identities
give
\[
   \lebesgue(g[B])
   =\mass\bigl(b^{-1}(g[B])\bigr)
   \ge\mass(F[D])\ge\mass(A)\ge1-\eps . \qedhere
\]
\end{proof}

\begin{thebibliography}{HPZZ}
\bibitem[DH]{DH}
R.~G. Downey and D.~R. Hirschfeldt,
\emph{Algorithmic Randomness and Complexity}, Theory and Applications of
Computability, Springer, New York, 2010.
\href{https://doi.org/10.1007/978-0-387-68441-3}{doi:10.1007/978-0-387-68441-3}.
\bibitem[HPZZ]{HPZZ}
P. Holick\'y, S.~P. Ponomarev, L. Zaj\'i\v cek and M. Zelen\'y,
Structure of the set of continuous functions with Luzin's property (N),
\emph{Real Anal. Exchange} \textbf{24} (1998/99), no.~2, 635--656.

\bibitem[PWY]{PWY}
A. Pauly, L. Westrick and L. Yu,
Luzin's (N) and randomness reflection,
\emph{J. Symb. Log.} \textbf{87} (2022), no.~2, 802--828.
\href{https://doi.org/10.1017/jsl.2020.73}{doi:10.1017/jsl.2020.73}.
\bibitem[W]{W}
L. Westrick,
\emph{Luzin's (N) and randomness reflection},
Midwest Computability Seminar with the CTA Online Seminar,
December~8, 2020, slides~9 and~15.
\href{https://math.uchicago.edu/~drh/mcs/westrick_slides.pdf}{Seminar slides}.
\end{thebibliography}
\end{document}
